\section{Introduction}
  \label{sec:introduction}

  Superconductivity remains one of the most striking collective phenomena in condensed matter physics.
  Despite more than a century of experimental and theoretical advances,
  a unified structural understanding across conventional and strongly correlated regimes
  remains elusive.
  In conventional low-$T_c$ materials, Bardeen--Cooper--Schrieffer (BCS) theory
  successfully explains pairing as a phonon-mediated instability of the Fermi surface.
  In contrast, high-$T_c$ superconductors such as cuprates and related materials
  exhibit strong electronic correlations, pseudogap behavior,
  and unconventional gap symmetries that resist a purely mediator-based description.

  The conceptual separation between these regimes has led to
  a proliferation of mechanisms tailored to specific material families.
  Phonon exchange, spin fluctuations, orbital fluctuations,
  and various strong-coupling approaches
  have been proposed as primary drivers of pairing.
  While many of these models capture important aspects of experimental data,
  they typically operate within a momentum-space interaction framework
  and often rely on material-specific assumptions.

  In this work we pursue a different route.
  We propose that superconductivity can be understood as
  \emph{projective phase locking}
  arising from a real-space geometric constraint.
  Within the Cosmochrony framework,
  electromagnetism is identified with the $U(1)$ fiber
  of a projective structure $\Pi \cong S^3$,
  whose phase coordinate $\theta$
  encodes the effective gauge degree of freedom.
  Topologically non-trivial excitations of this fiber
  correspond to charged solitonic configurations.
  A bound composite of winding number $w=2$
  emerges as a structurally distinct class,
  carrying effective charge $2e$ and supporting
  a collective phase degree of freedom.

  The central claim of this paper is that superconductivity
  results from the stabilization and phase locking
  of such $w=2$ composites.
  In conventional materials,
  a bosonic mediator such as a phonon
  lowers the projective gap $\Delta_\Pi$
  and stabilizes the composite configuration.
  In strongly correlated systems,
  pairing instead arises as a reduction of
  \emph{real-space frustration} in the admissibility
  constraints of the fiber.
  In both regimes,
  the London and Ginzburg--Landau structures
  follow from the cost of breaking global fiber coherence,
  rather than from an assumed symmetry-breaking potential.

  This approach yields a unified structural picture with several consequences.
  First, it naturally separates amplitude and phase scales:
  the projective gap governing local composite stability
  can differ from the global phase-locking temperature,
  providing a geometric account of pseudogap behavior.
  Second, the symmetry of the superconducting gap
  is determined by a real-space minimization problem
  on the fiber raccordement cost under lattice constraints,
  rather than by a specific momentum-dependent interaction kernel.
  Third, the phase stiffness scales with the density of active composites,
  leading to $\rho_s \propto \delta J$ in strongly correlated quasi-two-dimensional systems,
  where $\delta$ is the carrier concentration
  and $J$ the dominant frustration scale.

  A central predictive application of this framework
  concerns superconducting nickelates.
  Because their antiferromagnetic frustration sector
  is reduced and partially isotropized relative to cuprates,
  the fiber-raccordement criterion predicts
  an $s^\pm$ gap symmetry rather than
  $d_{x^2-y^2}$ symmetry at optimal doping.
  The predicted $T_c$ ratio between nickelates and cuprates
  follows from independently measurable parameters,
  including magnetic exchange scales and structure factors,
  without introducing free fitting constants.
  The predicted intermediate sensitivity to non-magnetic disorder
  provides a directly testable signature.

  The purpose of this paper is therefore threefold.
  We first derive the universal London and Ginzburg--Landau structure
  from the projective fiber geometry.
  We then extend the mechanism to the strongly correlated regime
  through a real-space frustration criterion.
  Finally, we formulate explicit falsifiable predictions,
  most notably for nickelate superconductors,
  and delineate the experimental conditions under which
  the mechanism would be ruled out.

  The framework presented here does not attempt
  to replace microscopic electronic models.
  Rather, it proposes a geometric organizing principle
  that constrains the admissible pairing channels
  and the scaling of phase coherence.
  Its validity can therefore be assessed
  by the consistency of its structural predictions
  with symmetry selection, stiffness scaling,
  disorder response, and pressure dependence
  across material families.
